\documentclass[12pt]{article}
\usepackage{amsmath, amssymb}
\usepackage[margin=1in]{geometry}
\setlength{\parskip}{6pt}
\title{QUBO Formulation: Robot Arm Path Planning}
\date{}
\begin{document}
\maketitle

\subsection*{Robot Arm Path Planning}

\paragraph{Problem Statement.}
Given a discrete time horizon $T$ and a sequence of state sets $\mathcal{S}_t$ for each time step $t \in \{0, \dots, T-1\}$, the Robot Arm Path Planning problem seeks a sequence of states $(s_0, s_1, \dots, s_{T-1})$ with $s_t \in \mathcal{S}_t$ that minimizes the sum of state-dependent obstacle costs and transition-dependent motion costs, subject to validity constraints on consecutive state pairs and fixed initial and final states.

\paragraph{Decision Variables.}
For each time step $t \in \{0, \dots, T-1\}$ and state $s \in \mathcal{S}_t$, introduce a binary variable $x_{t,s} \in \{0,1\}$. The variable $x_{t,s}$ equals $1$ if and only if state $s$ is selected at time step $t$, and $0$ otherwise.

\paragraph{QUBO Derivation.}
\textbf{Step 1 --- Original Objective.} The problem is to minimize the total obstacle cost plus the total motion cost:
\begin{align}
\min \quad & \sum_{t=0}^{T-1} \sum_{s \in \mathcal{S}_t} C_{\text{obs}}(s) x_{t,s} + \sum_{t=0}^{T-2} \sum_{s \in \mathcal{S}_t} \sum_{s' \in \mathcal{S}_{t+1}} C_{\text{mot}}(s,s') x_{t,s} x_{t+1,s'}.
\end{align}

\textbf{Step 2 --- Constraints.} The feasible region is defined by:
\begin{align}
\text{(C1)} \quad & \sum_{s \in \mathcal{S}_t} x_{t,s} = 1, \quad \forall t \in \{0, \dots, T-1\}, \\
\text{(C2)} \quad & x_{t,s} + x_{t+1,s'} \le 1, \quad \forall \text{ invalid transitions } (s,s') \text{ between } t \text{ and } t+1, \\
\text{(C3)} \quad & x_{0, s_{\text{start}}} = 1, \\
\text{(C4)} \quad & x_{T-1, s_{\text{goal}}} = 1.
\end{align}

\textbf{Step 3 --- Penalty Conversion.} We convert each constraint into a quadratic penalty term added to the objective.
For the equality constraint (C1), we add $\lambda_{\text{eq}} \left( \sum_{s \in \mathcal{S}_t} x_{t,s} - 1 \right)^2$. Using the idempotent property $x_{t,s}^2 = x_{t,s}$ for binary variables, we expand:
\begin{align}
\left( \sum_{s \in \mathcal{S}_t} x_{t,s} - 1 \right)^2 &= \sum_{s \in \mathcal{S}_t} x_{t,s}^2 + 2 \sum_{\substack{s, s' \in \mathcal{S}_t \\ s < s'}} x_{t,s} x_{t,s'} - 2 \sum_{s \in \mathcal{S}_t} x_{t,s} + 1 \nonumber \\
&= -\sum_{s \in \mathcal{S}_t} x_{t,s} + 2 \sum_{\substack{s, s' \in \mathcal{S}_t \\ s < s'}} x_{t,s} x_{t,s'} + 1.
\end{align}
For the inequality constraint (C2), since $x_{t,s}, x_{t+1,s'} \in \{0,1\}$, the condition $x_{t,s} + x_{t+1,s'} \le 1$ is violated if and only if both variables equal $1$. Thus, the penalty is simply $\lambda_{\text{ineq}} x_{t,s} x_{t+1,s'}$.
For the fixed start and goal constraints (C3) and (C4), we use $\lambda_{\text{start}} (x_{0, s_{\text{start}}} - 1)^2 = \lambda_{\text{start}} (-x_{0, s_{\text{start}}} + 1)$ and $\lambda_{\text{goal}} (x_{T-1, s_{\text{goal}}} - 1)^2 = \lambda_{\text{goal}} (-x_{T-1, s_{\text{goal}}} + 1)$.

\textbf{Step 4 --- Combined Hamiltonian.} Summing the objective and all penalty contributions yields the QUBO Hamiltonian $H(\mathbf{x})$:
\begin{align}
H(\mathbf{x}) &= \sum_{t=0}^{T-1} \sum_{s \in \mathcal{S}_t} C_{\text{obs}}(s) x_{t,s} + \sum_{t=0}^{T-2} \sum_{s \in \mathcal{S}_t} \sum_{s' \in \mathcal{S}_{t+1}} C_{\text{mot}}(s,s') x_{t,s} x_{t+1,s'} \nonumber \\
&\quad + \sum_{t=0}^{T-1} \lambda_{\text{eq}} \left( -\sum_{s \in \mathcal{S}_t} x_{t,s} + 2 \sum_{\substack{s, s' \in \mathcal{S}_t \\ s < s'}} x_{t,s} x_{t,s'} + 1 \right) \nonumber \\
&\quad + \sum_{\text{invalid } (t,s,s')} \lambda_{\text{ineq}} x_{t,s} x_{t+1,s'} + \lambda_{\text{start}} (-x_{0, s_{\text{start}}} + 1) + \lambda_{\text{goal}} (-x_{T-1, s_{\text{goal}}} + 1).
\end{align}

\textbf{Step 5 --- Q Matrix Entries and Offset.} Collecting linear and quadratic coefficients gives the general closed-form expressions for the upper-triangular Q matrix and constant offset $c$:
\begin{align}
Q_{(t,s),(t,s)} &= C_{\text{obs}}(s) - \lambda_{\text{eq}} + \lambda_{\text{start}} \delta_{t,0} \delta_{s, s_{\text{start}}} + \lambda_{\text{goal}} \delta_{t,T-1} \delta_{s, s_{\text{goal}}}, \\
Q_{(t,s),(t+1,s')} &= C_{\text{mot}}(s,s') + \lambda_{\text{ineq}} \cdot \mathbb{I}((s,s') \text{ is invalid}), \\
Q_{(t,s),(t,s')} &= 2\lambda_{\text{eq}} \quad \text{for } s \neq s', \\
Q_{(t,s),(t',s')} &= 0 \quad \text{otherwise}, \\
c &= T\lambda_{\text{eq}} + \lambda_{\text{start}} + \lambda_{\text{goal}}.
\end{align}
The QUBO is then formulated as $\min_{\mathbf{x} \in \{0,1\}^n} \mathbf{x}^\top Q \mathbf{x} + c$.

\paragraph{Penalty Weight Justification.}
The penalty weights must dominate any potential reduction in the original objective function to enforce feasibility. Specifically, $\lambda_{\text{eq}} > \max_{s} C_{\text{obs}}(s) \cdot \max_{t} |\mathcal{S}_t|$ ensures that violating the one-state-per-step constraint costs more than the maximum possible decrease in obstacle costs. $\lambda_{\text{ineq}} > \max_{s,s'} C_{\text{mot}}(s,s')$ guarantees that the penalty for an invalid transition exceeds any motion cost savings. Finally, $\lambda_{\text{start}}, \lambda_{\text{goal}} > \max_{s} C_{\text{obs}}(s)$ ensures the start and goal states are strictly fixed.

\paragraph{Correctness Argument.}
Minimizing $H(\mathbf{x})$ over $\{0,1\}^n$ recovers the optimal path because any feasible solution satisfies all constraints exactly, yielding zero penalty energy. Any infeasible solution incurs a penalty strictly greater than the maximum possible decrease in the original objective function (by the choice of $\lambda$ parameters). Consequently, the global minimum of $H(\mathbf{x})$ cannot be infeasible; it must correspond to a feasible solution that minimizes the original cost function.

\paragraph{Illustrative Example.}
Consider a concrete instance with $T=3$ time steps and $|\mathcal{S}_t|=3$ states per step, labeled $\{0,1,2\}$. Variables are flattened to a single index $k = 3t + s$, yielding $n=9$ variables $x_0, \dots, x_8$. Assume uniform obstacle costs $C_{\text{obs}}(s) = c_o$ and uniform motion costs $C_{\text{mot}}(s,s') = c_m$. Suppose the only forbidden transition is $(s=0, s'=0)$ between $t=0$ and $t=1$, corresponding to variables $x_0$ and $x_3$. The start state is $s_{\text{start}}=0$ ($x_0$) and the goal state is $s_{\text{goal}}=2$ ($x_8$).

The concrete objective is:
\begin{align}
\mathcal{O}(\mathbf{x}) &= c_o \sum_{k=0}^{8} x_k + c_m (x_0x_3 + x_0x_4 + x_0x_5 + x_1x_3 + x_1x_4 + x_1x_5 + x_2x_4 + x_2x_5 + x_3x_6 + x_3x_7 + x_4x_6 + x_4x_7 + x_4x_8 + x_5x_7 + x_5x_8).
\end{align}
The concrete penalty terms expand to:
\begin{align}
\mathcal{P}_{\text{eq}} &= \lambda_{\text{eq}} \left( -\sum_{k=0}^{8} x_k + 2(x_0x_1 + x_0x_2 + x_1x_2 + x_3x_4 + x_3x_5 + x_4x_5 + x_6x_7 + x_6x_8 + x_7x_8) + 3 \right), \\
\mathcal{P}_{\text{ineq}} &= \lambda_{\text{ineq}} x_0 x_3, \\
\mathcal{P}_{\text{start}} &= \lambda_{\text{start}} (-x_0 + 1), \quad \mathcal{P}_{\text{goal}} = \lambda_{\text{goal}} (-x_8 + 1).
\end{align}
Combining these, the full Hamiltonian $H(\mathbf{x}) = \mathcal{O}(\mathbf{x}) + \mathcal{P}_{\text{eq}} + \mathcal{P}_{\text{ineq}} + \mathcal{P}_{\text{start}} + \mathcal{P}_{\text{goal}}$ collects coefficients for each monomial. The optimal feasible path selects states $s_0=0, s_1=1, s_2=2$, corresponding to the binary vector $\mathbf{x}^* = (1,0,0, 0,1,0, 0,0,1)$. Substituting $\mathbf{x}^*$ into $H(\mathbf{x})$ yields the minimum energy:
\begin{align}
H(\mathbf{x}^*) = 3c_o + 2c_m.
\end{align}
This matches the sum of obstacle costs for three steps and motion costs for two transitions, confirming the formulation correctly encodes the instance.

\end{document}